\chapter*{Abstract}
\addcontentsline{toc}{chapter}{Abstract}

Distributed Internet of Things (IoT) sensor networks frequently operate under tight energy
constraints, particularly in large-scale smart building deployments. This thesis proposes an
adaptive, machine-learning-based scheduling framework that dynamically adjusts sensor sampling
rates in response to network load patterns and ambient environmental variance. The framework is
evaluated on a real-world dataset collected from a Nexa Systems smart building deployment in
Kuala Lumpur, comprising 48 sensors over a 90-day observation window. Experimental results show a
27.4\% reduction in aggregate energy consumption compared to conventional fixed-rate scheduling,
with a negligible increase in average data latency (under 120 milliseconds). The key contributions
of this work are: (i) a lightweight scheduling algorithm suitable for deployment on
resource-constrained edge gateways, and (ii) a reusable evaluation methodology applicable to other
sensor network topologies.

\medskip
\noindent\textbf{Keywords:} Internet of Things, machine learning, energy efficiency, edge
computing, adaptive scheduling

% ---------- Acknowledgements ----------
